\documentclass[10pt,twocolumn]{article}

\usepackage[margin=0.78in,columnsep=0.28in]{geometry}
\usepackage{times}
\usepackage{amsmath}
\usepackage{graphicx}
\usepackage{booktabs}
\usepackage{multirow}
\usepackage{hyperref}
\usepackage[T1]{fontenc}
\usepackage{titlesec}
\usepackage{caption}
\usepackage{authblk}
\usepackage{enumitem}
\usepackage{placeins}

\renewcommand{\dbltopfraction}{0.95}
\renewcommand{\dblfloatpagefraction}{0.8}
\setcounter{dbltopnumber}{3}
\titlespacing*{\section}{0pt}{8pt}{4pt}
\titlespacing*{\subsection}{0pt}{6pt}{3pt}
\captionsetup{font=small,labelfont=bf}

\title{\vspace{-1.2em}Does Fairness Improve RAG Performance, or Is It Diversity?\\
\large A Critical Reproduction and Extension of \emph{Towards Fair RAG}\vspace{-0.3em}}

\author[1]{Idan Peretz}
\author[1]{Avi Simkin}
\affil[1]{\normalsize \{308560150, 312485816\}}
\date{}

\begin{document}
\maketitle
\vspace{-2.2em}

\begin{abstract}
\noindent
Fairness-aware stochastic rerankers are usually expected to trade ranking quality for
more equitable item exposure, yet Eun Kim and Diaz (2025) report that RAG systems
built on such rerankers can maintain, and sometimes exceed, the generation utility of
deterministic ranking. We reproduce this setup on the LaMP benchmark and investigate
the mechanism behind it, contrasting a deterministic diversity-relevance re-ranker
(MMR) against pure stochastic sampling (Plackett-Luce) to separate the effect of
diversity from the effect of randomness itself. We find that relevance-aware
diversification is a real, though narrow and task-dependent, source of utility gain
over deterministic ranking, while blind randomization mostly is not --- evidence
against a simple story in which any diversification helps equally. Building on this,
we propose RandMMR, a hybrid reranker combining stochastic sampling with
relevance-aware diversification, and find it pushes fairness substantially further
than either ingredient alone without a significant utility cost, and with an
aggregate utility picture close to neutral --- no cell loses significantly, two
win significantly, both on the same (generator, ranker) pair. We also offer a
critical reading of
both the
original paper's evaluation protocol and of the LaMP benchmark itself as ground truth
for this kind of study, and lay out concrete directions for future work.
\end{abstract}

\vspace{0.3em}

\section{Introduction and Problem Formulation}
\label{sec:intro}

Fair-ranking research typically formalizes an unavoidable trade-off: enforcing more
equitable exposure among relevant items costs some ranking quality relative to a
purely relevance-ordered list \cite{castillo2018fairness,diaz2020evaluating}.
Eun Kim and Diaz \cite{eunkim2025fairrag} extend this literature to
retrieval-augmented generation (RAG), reporting that RAG systems using a stochastic,
fairness-aware retriever policy can maintain, and in some regimes even exceed, the
generation utility of a fully deterministic retriever, despite that trade-off. This
motivates two questions, the first stated in our original project proposal:

\begin{itemize}[leftmargin=1.1em]
  \item \textbf{RQ1 (Mechanism).} If fairness interventions do not cost utility, is
  this because they increase the \emph{diversity} of retrieved items --- intuitively,
  a generator given a less redundant document set has more distinct evidence to draw
  on?
  \item \textbf{RQ2 (Improvement).} Given whatever RQ1 reveals about \emph{why}
  fairness interventions do not cost utility, can that insight be used to design a
  reranking policy that pushes fairness (lower EE-D) further than the settings we
  study reach on their own, without giving utility back up in exchange?
\end{itemize}

Formally, following \cite{eunkim2025fairrag}, a RAG system is a tuple
$\langle R, G, k \rangle$ of a retriever $R$, generator $G$, and truncation depth
$k$. A deterministic $R$ returns one ranked list per query; a stochastic policy $S$
instead samples a ranking $\pi \sim S(\cdot \mid q)$, and system quality is reported
in expectation over $N$ sampled rankings: Expected Exposure Disparity/Relevance
(EE-D/EE-R, item-side ranking fairness and quality) and Expected Utility (EU,
downstream generation quality), each normalized to a theoretical or empirical bound
per query \cite{diaz2020evaluating,eunkim2025fairrag}. We add one further quantity
not defined in the source paper: \emph{intra-list diversity}
$\mathrm{ILD}(\pi) = 1 - \overline{\mathrm{Jaccard}}(\pi)$, the mean pairwise
token-set dissimilarity of the $k$ retrieved items, following the standard
similarity-vs-diversity framing of \cite{smyth2001similarity}. RQ1 asks whether EU is
better predicted by ILD than by the fairness dial itself.

\section{Related Work}
\label{sec:related}

\textbf{Fair and stochastic ranking.} Fairness in ranking is not simply fairness in
classification applied item-by-item: a ranked list's effect on an item is mediated by
\emph{exposure}, computed via the position-weighted formula defined by Diaz et al.\
\cite{diaz2020evaluating} rather than by an item's bare relevance. Item-side exposure
fairness is formalized via the Expected Exposure framework \cite{diaz2020evaluating}, which
measures a policy's exposure allocation against a fair target using a stochastic
ranking sampled by the Plackett-Luce (PL) model
\cite{plackett1975analysis,luce1959individual}; the temperature parameter $\alpha$
that we sweep here controls how close PL sampling is to the deterministic base
ranking. Castillo \cite{castillo2018fairness} surveys the fairness-transparency
trade-off in rankings that motivates this line of work. \emph{Towards Fair RAG}
\cite{eunkim2025fairrag} is the paper under reproduction: it applies this fairness
framework to RAG, measuring EE-D/EE-R/EU across LaMP tasks and several
retriever-generator pairs, and is the direct source of our experimental design,
metrics, and normalization convention (Appendix C of \cite{eunkim2025fairrag}). Its
utility metric, EU, is computed per sampled ranking as the downstream generator's
task-appropriate score (accuracy, MAE, or ROUGE-L, depending on the LaMP task) when
conditioned on that ranking's retrieved items, then normalized per query against an
oracle-ranking upper bound before being averaged across $N$ sampled rankings.

\textbf{Diversity in retrieval.} Maximal Marginal Relevance (MMR)
\cite{carbonell1998use} is the classic deterministic diversity-relevance trade-off
re-ranker, greedily selecting each next item to maximize $\lambda \cdot
\mathrm{relevance} - (1-\lambda)\cdot\max\text{-similarity-to-selected}$; we use it
as a second, non-stochastic diversification arm alongside PL. Intra-list diversity
via pairwise dissimilarity follows \cite{smyth2001similarity}.

\textbf{Personalization benchmark and models.} We use LaMP
\cite{salemi2023lamp} as our dataset (Section~\ref{sec:data}) and Flan-T5
\cite{chung2024scaling} as our generator family, matching \cite{eunkim2025fairrag}.

\section{Dataset}
\label{sec:data}

We use LaMP \cite{salemi2023lamp}, the same benchmark as \cite{eunkim2025fairrag},
for two reasons. First, and primarily, paper-faithfulness: LaMP is the only dataset
for which the original paper published per-task, per-generator utility-labeled
corpora (filtered to items with a measurable generator-specific utility delta), so
reproducing on it is the only way to compare our numbers directly against theirs
rather than against a different task distribution. Second, tractability: LaMP's
candidate pools are small enough that deriving actual relevance labels is
computationally feasible in the first place (Section~\ref{sec:threats}), which gives
us, for the same queries, both a retrieval-fairness signal (which documents are
labeled useful) and a generation-correctness signal --- a joint
retrieval-\emph{and}-generation labeling that,
to our knowledge, is not available together in other RAG-evaluation datasets, and is
exactly what studying fairness in RAG needs. LaMP additionally spans seven
personalization tasks with heterogeneous metrics (accuracy, MAE, ROUGE-L), letting us
check whether any finding is task-specific or general.

We verified our local corpus against the original paper's own published statistics
(Appendix D, Tables 4--5): per-task query count, average candidate-document count,
average useful-document count, and average useful-document percentage all match
to the paper's reported precision across all 7 tasks $\times$ 2 generators (28
cells), confirming our filtered dataset is the identical one used upstream, not a
re-derived approximation.

We run two generators (Flan-T5-Small, Flan-T5-Base \cite{chung2024scaling}) and two
retrievers (BM25, Contriever), against the paper's larger grid (adding SPLADE and
two larger generators) --- a scope reduction driven by available compute, not
methodology. Reranking settings are deterministic, a gold/oracle ranker, PL at
$\alpha \in \{1,2,4,8\}$, and MMR at $\lambda \in \{0.3,0.45,0.55,0.7,0.85\}$ (a sixth
setting, $\lambda{=}0.15$, was also run but excluded from analysis as too
diversity-heavy to be informative at the low end the tests below turn on), all
truncated to $k=5$, matching \cite{eunkim2025fairrag}'s reported $k$.

\subsection{Limitations of the LaMP Ground Truth}
\label{sec:threats}

\begin{table*}[t]
\centering
\scriptsize
\caption{LaMP candidate-pool size and useful- (i.e.\ relevant-) document rate, per task and generator (from each task's relevance-mapping file).}
\label{tab:corpus}
\begin{tabular}{@{}lrrrrrrrrrrrrrr@{}}
\toprule
& \multicolumn{2}{c}{LaMP-1} & \multicolumn{2}{c}{LaMP-2} & \multicolumn{2}{c}{LaMP-3} & \multicolumn{2}{c}{LaMP-4} & \multicolumn{2}{c}{LaMP-5} & \multicolumn{2}{c}{LaMP-6} & \multicolumn{2}{c}{LaMP-7} \\
Generator & cand. & \%use & cand. & \%use & cand. & \%use & cand. & \%use & cand. & \%use & cand. & \%use & cand. & \%use \\
\midrule
Flan-T5-Small & 123.5 & 9.5 & 52.8 & 22.5 & 189.8 & 34.3 & 192.2 & 27.1 & 106.1 & 24.8 & 86.0 & 35.9 & 19.4 & 45.5 \\
Flan-T5-Base  & 102.9 & 22.5 & 45.6 & 29.9 & 185.3 & 41.2 & 187.0 & 31.6 & 105.6 & 25.7 & 86.2 & 38.7 & 21.7 & 33.0 \\
\bottomrule
\end{tabular}
\end{table*}

These reasons for choosing LaMP come with limitations worth stating up front, shared
with the source paper and not introduced by our reproduction, since they bound how
far any effect we measure in the results that follow (Section~\ref{sec:results}),
including RQ1's diversity test and RQ2's RandMMR design alike, should generalize.

\textbf{Relevance is a per-generator utility delta, not a fixed judgment.} Unlike
TREC-style collections, where relevance is a human judgment shared by every system, a
LaMP document is labeled relevant iff including it, versus a zero-context baseline,
\emph{improves that specific generator's} task score
($\delta = \text{augment\_score} - \text{baseline\_score}$; label $=1$ iff
$\delta{>}0$). Ground truth is therefore generator-specific: on the 21 LaMP-1 queries
shared by both generators' filtered sets, 22.8\% of document labels flip between
them, and the useful-document rate itself differs sharply (9.4\% vs.\ 20.2\%). Any
fairness-ranking result computed against these labels is a statement about one
model's learned usefulness, not a model-agnostic ground truth.

\textbf{The corpus is small, with an unusually high useful-document rate.} Table
\ref{tab:corpus} gives per-task, per-generator statistics: candidate pools of
19--192 documents/query (mean $\sim$110) and useful-document rates of
9.5\%--45.5\% (mean $\sim$29\%), against classical TREC collections' orders-of-
magnitude-larger pools and sub-1\% relevance rates. At $k{=}5$
\cite{eunkim2025fairrag} and a $\sim$30\% useful rate, even a fairness-motivated
rerank that displaces much of the top-$k$ still has decent odds of retaining several
useful documents, narrowing the deterministic-vs-stochastic utility gap the whole
fairness/utility trade-off is measured against.

\textbf{Our own generators are small; the source paper's are dated.} Flan-T5-Small
(77M) and Flan-T5-Base (250M) \cite{chung2024scaling}, the two we run, are one to two
orders of magnitude smaller than production RAG generators ($\geq$3B parameters).
The source paper's own larger generators are a separate limitation, not the same
one: its largest, Flan-UL2 ($\sim$20B), is not undersized so much as outdated ---
the same 2022-era family, not finetuned beyond LaMP's prompt format, and several
generations behind what is typically deployed today. Whether either paper's findings
persist at production scale is untested by both (see Future Work).

\textbf{The EE-D/EE-R trade-off is not a per-query guarantee.} The same small,
high-relevance corpus lets PL's aggregate trade-off fail query-by-query while holding
on average: an individual query's PL EE-R exceeds its own deterministic EE-R in
29.7\% of observations (37--38\% at $\alpha\in\{1,2,4\}$, still 8.3\% at the mildest
$\alpha{=}8$) --- though this mostly washes out once pooled to a macro mean (14 of
112 conditions, all at $\alpha{=}8$, $\Delta \leq 0.004$). A caution against
over-interpreting any single query's or run's numbers.

\textbf{ROUGE-L is an imperfect stand-in for generation quality.} For the four
text-generation tasks (LaMP-4--7), both our utility labels and EU itself score
correctness via ROUGE-L, an $n$-gram-overlap metric --- it can score a valid
paraphrase of the gold reference lower than a wrong but lexically similar answer,
rewarding surface overlap over semantic correctness. Errors from this mismatch
propagate into the ground truth itself for these four tasks, not just into our
reported EU. A more faithful metric, an LLM-as-judge, exists but was computationally
out of reach at this project's scale; ROUGE-L is the standard LaMP-benchmark metric
for these tasks \cite{salemi2023lamp}, not a choice specific to either paper.

\section{Methodology Notes}
\label{sec:method}

We treat the original paper's open-sourced code as a direct reference throughout:
matching its library versions, reusing its utility-labeling technique
(Section~\ref{sec:threats}) and implementation details (prompt templates, scoring
functions) wherever our framework's structure allowed it directly, and otherwise
re-implementing to match its stated method as closely as we could. Reproduction
fidelity is imperfect on the two tasks for which the source paper publishes
deterministic-baseline EU directly, driven mainly by Flan-T5-Small's own
instability rather than by task --- a content-invariant failure mode on LaMP-1
(it predicts the same label regardless of retrieved context, which also
destabilizes its normalization ceiling) and, on LaMP-4, a gap between our
normalization ceiling and the paper's fuller one. We treat both as limits of
exact reproduction, not codebase errors; the full per-cell comparison against
the paper's own numbers is in Table~\ref{tab:baseline}, Appendix~\ref{app:baseline}.

Every EU$_{\text{norm}}$ value in this paper uses the source paper's own per-query
formula, $\text{EU}_{\text{norm}}(q){=}\text{EU}(q)/\max(\text{gold\_ceiling}(q),
\text{model\_ceiling}(q))$ (Appendix C of \cite{eunkim2025fairrag}): raw utility
divided by the higher of a gold-ranked ceiling (3 PL draws over the oracle ranking
at $\alpha{=}8$, a lower bound on the paper's own fuller stochastic search) and a
model ceiling (the best output any rerank method we ran ever produced for that
query, at that generator and ranker --- PL at any $\alpha$, MMR at any $\lambda$,
RandMMR, or deterministic, not only PL: Appendix C of \cite{eunkim2025fairrag}
defines the bound for a fixed $\text{RAG}(R,G,k)$, and every method we run at a
given $(R,G,k)$ is a real empirical sample of what that same system can produce),
falling back to $1.0$ when both ceilings are zero.

For the diversity-mediation test (RQ1), for each LaMP task we bin PL's and MMR's own
generated lists \emph{separately} into ILD$_{\text{norm}}$ quintiles and compare the
least-diversified quintile's (Q1) mean EU$_{\text{norm}}$ against the deterministic
baseline's, via a two-sided Welch's $t$-test, task-stratified with equal weight across
all four generator$\times$retriever cells. Because a query can land in the same bin
under more than one $\alpha$ or $\lambda$ setting, these per-bin tests treat repeated
observations of the same query as independent rather than applying a cluster-robust
correction, so reported significance is a mild overstatement relative to a fully
query-clustered test. One MMR setting, $\lambda{=}1.0$, is excluded: at $\lambda{=}1$
MMR's diversity term has zero weight, so its greedy selection provably reproduces
the deterministic ranking exactly --- it is not a diversification condition at all. A second precision control matters
at this scale: some (generator, ranker, $\alpha$) conditions have been resampled more
times ($N{\in}\{20,30\}$) than others ($N{=}10$), and pooling each condition at its
own maximum $N$ independently would silently give some conditions tighter per-query
utility estimates than the ones they are compared against. Every PL-derived number in
Section~\ref{sec:results} instead pools all conditions inside one task at a single $N$
shared by every (generator, ranker, $\alpha$) triple for that task --- the highest
such $N$, or $N{=}10$ where no higher value is universal. This precision-matching
matters for EE-D too, not only EU: EE-D is computed jointly over all $N$ sampled
rankings for a query, so lower $N$ gives a noisier, clumpier empirical exposure
allocation that looks less fair and converges toward a fairer-looking value only
as $N$ grows; left uncontrolled, this would bias any $N$-mismatched comparison
toward making the higher-$N$ condition look artificially fairer for reasons
unrelated to the reranking policy itself.

\section{Results}
\label{sec:results}

\subsection{RQ1: Diversity as the mediating mechanism}
\label{sec:rq1}

RQ1 asks whether fairness's apparent lack of utility cost is really a diversity
effect in disguise. We first check the premise --- that PL's and MMR's dials actually
move measured diversity as intended --- before testing whether that diversity
predicts utility.

\begin{figure}[t]
\centering
\includegraphics[width=\linewidth]{fig_diversity_lamp4_panel.png}
\caption{ILD-Jaccard diversity by rerank setting, LaMP-4 (representative; all seven
tasks show the same qualitative pattern, largest swing on LaMP-6 and smallest on
LaMP-7). Diamond = mean, black line = median. PL/MMR colors match Figure~\ref{fig:mmr-vs-pl-q1}'s legend.}
\label{fig:diversity}
\end{figure}

Figure~\ref{fig:diversity} shows LaMP-4 as a representative example: PL and MMR both
raise list diversity as their dial is relaxed ($\alpha{:}8{\to}1$,
$\lambda{:}0.85{\to}0.3$). The other six tasks show the same qualitative pattern, but
the swing's size is task-dependent --- largest for LaMP-6 (mean ILD $0.65 \to 0.82$)
and smallest for LaMP-7 (everything within $0.92$--$0.95$), consistent with LaMP-7
having very few candidate documents per query ($\sim$20, Table~\ref{tab:corpus}) and
thus little room to diversify.

\begin{table*}[t]
\centering
\scriptsize
\caption{PL and MMR utility gain over deterministic, by each generator/ranker/method
cell's own ILD$_{\text{norm}}$ quartile, LaMP-1 and LaMP-4 (the two tasks with a
published baseline). Bracketed values are the quartile's own ILD$_{\text{norm}}$
boundaries. $^*p{<}0.05$, $^{**}p{<}0.01$ (Welch's $t$-test vs.\ deterministic); the
other five LaMP tasks show a similar pattern, given in full in
Appendix~\ref{app:ild-quartiles-full} (Table~\ref{tab:ild-quartiles-full}).}
\label{tab:ild-quartiles}
\begin{tabular}{@{}lllrlrrrr@{}}
\toprule
Task & Gen. & Ranker & Base & Method & Q1 (least diverse) & Q2 & Q3 & Q4 (most diverse) \\
\midrule
\multirow{2}{*}{LaMP-1} & \multirow{2}{*}{T5-Base} & \multirow{2}{*}{BM25} & \multirow{2}{*}{0.681} & PL & [0.73,0.97] $-$0.071 & [0.97,0.98] $-$0.083$^{*}$ & [0.98,0.99] $-$0.121$^{**}$ & [0.99,1.00] $-$0.165$^{**}$ \\
 & & & & MMR & [0.73,0.97] $+$0.053 & [0.97,0.98] $-$0.039 & [0.98,0.99] $+$0.002 & [0.99,1.00] $-$0.089$^{*}$ \\
\cmidrule(lr){1-9}
\multirow{2}{*}{LaMP-1} & \multirow{2}{*}{T5-Base} & \multirow{2}{*}{Contr.} & \multirow{2}{*}{0.642} & PL & [0.81,0.96] $-$0.042 & [0.96,0.98] $-$0.061 & [0.98,0.99] $-$0.021 & [0.99,1.00] $-$0.081$^{*}$ \\
 & & & & MMR & [0.85,0.97] $-$0.007 & [0.97,0.98] $+$0.020 & [0.98,0.99] $+$0.067 & [0.99,1.00] $-$0.094$^{*}$ \\
\cmidrule(lr){1-9}
\multirow{2}{*}{LaMP-1} & \multirow{2}{*}{T5-Small} & \multirow{2}{*}{BM25} & \multirow{2}{*}{0.451} & PL & [0.75,0.96] $-$0.135 & [0.96,0.98] $+$0.051 & [0.98,0.99] $+$0.047 & [0.99,1.00] $-$0.071 \\
 & & & & MMR & [0.90,0.96] $-$0.143 & [0.96,0.98] $+$0.089 & [0.98,1.00] $+$0.002 & [1.00,1.00] $-$0.054 \\
\cmidrule(lr){1-9}
\multirow{2}{*}{LaMP-1} & \multirow{2}{*}{T5-Small} & \multirow{2}{*}{Contr.} & \multirow{2}{*}{0.431} & PL & [0.81,0.96] $-$0.043 & [0.96,0.97] $-$0.047 & [0.97,0.99] $-$0.022 & [0.99,1.00] $+$0.035 \\
 & & & & MMR & [0.87,0.96] $-$0.041 & [0.96,0.98] $+$0.031 & [0.98,0.99] $+$0.019 & [0.99,1.00] $+$0.022 \\
\midrule
\multirow{2}{*}{LaMP-4} & \multirow{2}{*}{T5-Base} & \multirow{2}{*}{BM25} & \multirow{2}{*}{0.291} & PL & [0.46,0.97] $-$0.011 & [0.97,0.98] $-$0.028$^{*}$ & [0.98,0.99] $-$0.035$^{**}$ & [0.99,1.00] $-$0.022 \\
 & & & & MMR & [0.46,0.97] $+$0.003 & [0.97,0.98] $+$0.021 & [0.98,0.99] $-$0.006 & [0.99,1.00] $-$0.005 \\
\cmidrule(lr){1-9}
\multirow{2}{*}{LaMP-4} & \multirow{2}{*}{T5-Base} & \multirow{2}{*}{Contr.} & \multirow{2}{*}{0.340} & PL & [0.50,0.97] $-$0.002 & [0.97,0.98] $-$0.060$^{**}$ & [0.98,0.99] $-$0.057$^{**}$ & [0.99,1.00] $-$0.071$^{**}$ \\
 & & & & MMR & [0.52,0.97] $+$0.029$^{*}$ & [0.97,0.99] $-$0.005 & [0.99,1.00] $-$0.033$^{*}$ & [1.00,1.00] $-$0.005 \\
\cmidrule(lr){1-9}
\multirow{2}{*}{LaMP-4} & \multirow{2}{*}{T5-Small} & \multirow{2}{*}{BM25} & \multirow{2}{*}{0.297} & PL & [0.47,0.97] $-$0.010 & [0.97,0.98] $-$0.021 & [0.98,0.99] $-$0.005 & [0.99,1.00] $-$0.001 \\
 & & & & MMR & [0.47,0.97] $+$0.006 & [0.97,0.98] $-$0.002 & [0.98,0.99] $-$0.009 & [0.99,1.00] $+$0.019 \\
\cmidrule(lr){1-9}
\multirow{2}{*}{LaMP-4} & \multirow{2}{*}{T5-Small} & \multirow{2}{*}{Contr.} & \multirow{2}{*}{0.326} & PL & [0.51,0.98] $-$0.004 & [0.98,0.99] $-$0.059$^{**}$ & [0.99,0.99] $-$0.033$^{**}$ & [0.99,1.00] $-$0.052$^{**}$ \\
 & & & & MMR & [0.53,0.97] $+$0.028 & [0.97,0.99] $-$0.012 & [0.99,1.00] $+$0.011 & [1.00,1.00] $-$0.047$^{**}$ \\
\bottomrule
\end{tabular}
\end{table*}

The key test is whether this diversity increase predicts \emph{higher} utility.
Table~\ref{tab:ild-quartiles}, inspired by the source paper's own Table 2 (which bins
by EE-D interval, pooled across all cells), instead bins PL's and MMR's own lists
into ILD$_{\text{norm}}$ quartiles, separately per method and, unlike the source
paper's pooled table, kept disaggregated by generator and ranker rather than pooled
across them --- reporting each quartile's actual boundary values alongside its
EU$_{\text{norm}}$ gain over deterministic, for the two tasks with a published
baseline (Table~\ref{tab:baseline}). PL's Q2--Q4 delta is negative in 21 of the 24
(task, cell, quartile) combinations --- every positive exception is on
Flan-T5-Small --- and reaches significance in 12: 9 of 12 on Flan-T5-Base, only 3
of 12 on the smaller, noisier Flan-T5-Small (Section~\ref{sec:threats}), and all
three of those from its one best-behaved cell (LaMP-4/Contriever). MMR's Q1 delta
is positive in 5 of the 8 cells, reaching significance in one
(Flan-T5-Base/Contriever, LaMP-4, $+0.029$, $p{<}0.05$) --- the same direction the
pooled quintile view in Figure~\ref{fig:mmr-vs-pl-q1} below shows, now visible cell
by cell rather than only in aggregate, even though most individual cells are too
small to clear $p{<}0.05$ on their own.

\begin{figure}[t]
\centering
\includegraphics[width=\linewidth]{fig_mmr_vs_pl_q1.pdf}
\caption{$\Delta$EU$_{\text{norm}}$ vs.\ deterministic at the Q1 (least-diversified)
bin, PL vs.\ MMR, per task. $^*p{<}0.05$, $^{**}p{<}0.01$.}
\label{fig:mmr-vs-pl-q1}
\end{figure}

Figure~\ref{fig:mmr-vs-pl-q1} breaks this down per task: MMR's Q1 delta vs.\
deterministic is positive in 5 of 7 tasks but significant in only one (LaMP-6,
$+0.039$, $p{<}0.01$); the other positive tasks are small and non-significant
(LaMP-3: $+0.008$; LaMP-4: $+0.012$; LaMP-5: $+0.006$; LaMP-7: $+0.020$), and the two
negative tasks are also non-significant (LaMP-1: $-0.022$; LaMP-2: $-0.028$). PL's Q1
delta is positive in only 2 of 7 tasks (LaMP-3, LaMP-7, both non-significant) and
negative in the other 5, reaching significance in one --- LaMP-6 ($-0.020$,
$p{<}0.05$); LaMP-1's negative delta ($-0.058$) falls just short of significance
($p{=}0.052$). LaMP-6 is the clearest single-task contrast in our results: the same
task where MMR's Q1 gain is significant is also the task where PL's Q1 change is a
significant \emph{loss}, relevance-aware and blind diversification pulling in
opposite directions on identical query sets.

This is a real but narrow finding, weaker than a broad confirmation of the
diversity-mediation hypothesis in our original proposal. The benefit, where it
exists, is concentrated in the least-diversified quartile/quintile (Q1) rather than
growing with more diversification, and it is sharply asymmetric between methods:
mild, relevance-aware diversification (MMR) can be a real source of utility gain,
but only in a minority of tasks, while PL's fairness-motivated randomization, on its
own, is if anything more likely to cost utility than to gain it.

\subsection{RQ2: RandMMR, a hybrid PL + randomized-MMR reranker}
\label{sec:rq2}

Motivated by the source paper's finding that fairness interventions do not reliably
cost utility, and by RQ1's finding that mild MMR-style diversification is a real (if
partial) reason why, we propose
\textbf{RandMMR}: PL sampling of the top-ranked item, then a randomized-MMR pass
($\lambda \sim \mathrm{Unif}(\ell_{lo},\ell_{hi})$) over the remaining ranks.
Randomizing only the top rank is deliberate, not a simplification: MMR selects each
later item by penalizing similarity to whatever is \emph{already} selected, so one
randomized draw at rank 1 cascades through every subsequent pick even though ranks
$2..k$ are themselves chosen deterministically given that draw --- in principle
trading less raw randomness for a comparable fairness gain, since the whole list's
exposure disparity moves with the top-rank draw, not just that one position's own
share of it. We
evaluate the RandMMR condition run across the most tasks and cells ($\alpha{=}4$,
$\lambda\sim\mathrm{Unif}(0.8,1.0)$, $N{=}20$: all 7 LaMP tasks, 26 completed
conditions --- the full 4 generator/retriever cells for six tasks, plus 2 of 4, both
Flan-T5-Base, for LaMP-6) against plain PL, choosing per
cell the PL $\alpha \in \{1,2,4,8\}$ whose own mean EE-D is closest to RandMMR's ---
biased toward an $\alpha$ with the same or slightly \emph{higher} disparity than
RandMMR's own, so PL never gets an unearned fairness advantage in the comparison ---
unless the single closest $\alpha$ is already within 0.02 EE-D, in which case that one
is used regardless of direction.

\begin{table}[t]
\centering
\scriptsize
\setlength{\tabcolsep}{3pt}
\caption{RandMMR ($\alpha{=}4$, $\lambda\sim U(0.8,1.0)$) vs.\ matched plain PL
($\alpha{=}2$ throughout; see text), LaMP-1 and LaMP-4 (macro means). EE-D and EU
are both normalized per query to $[0,1]$ (EE-D: the paper's own bound; EU: our
per-query empirical-max EU$_{\text{norm}}$). RandMMR's EE-D is lower (fairer) than
the matched PL's in every row shown. One of the eight $\Delta$EU values reaches
significance (Flan-T5-Small/BM25, LaMP-4, $p{<}0.01$, favoring RandMMR); the other
five LaMP tasks show a similar pattern, given in full in
Appendix~\ref{app:randmmr-full} (Table~\ref{tab:hybrid-full}).}
\label{tab:hybrid}
\begin{tabular}{@{}ll rr r@{}}
\toprule
& & \multicolumn{2}{c}{EE-D/EU} & \\
\cmidrule(lr){3-4}
Task & Gen./Rank. & PL & RandMMR & $\Delta$EU \\
\midrule
\multirow{4}{*}{LaMP-1} & Base/BM25    & 0.19/0.51 & 0.19/0.49 & $-$0.019 \\
 & Base/Contr.  & 0.18/0.56 & 0.17/0.51 & $-$0.048 \\
 & Small/BM25   & 0.18/0.40 & 0.17/0.39 & $-$0.008 \\
 & Small/Contr. & 0.17/0.39 & 0.16/0.41 & $+$0.020 \\
\cmidrule(lr){1-5}
\multirow{4}{*}{LaMP-4} & Base/BM25    & 0.25/0.25 & 0.25/0.25 & $+$0.003 \\
 & Base/Contr.  & 0.24/0.25 & 0.24/0.25 & $+$0.000 \\
 & Small/BM25   & 0.26/0.22 & 0.25/0.27 & $+$0.050$^{**}$ \\
 & Small/Contr. & 0.24/0.26 & 0.23/0.26 & $-$0.007 \\
\bottomrule
\end{tabular}
\end{table}

Table~\ref{tab:hybrid} details the two tasks with a published baseline; the matched
$\alpha$ is 2 for all eight rows (see text above for the matching rule). RandMMR's
EU$_{\text{norm}}$ exceeds the matched PL's in 4 of these 8 rows, one significantly
(Flan-T5-Small/BM25, LaMP-4, $+0.050$, $p{<}0.01$). Across the remaining 20 rows
(LaMP-2, 3, 5, 6, 7), it does so in 8 of 20, one more significantly
(Flan-T5-Small/BM25, LaMP-2, $+0.108$, $p{<}0.05$) --- 12 of 28 completed cells
overall (43\%), both significant cells still favoring RandMMR and none favoring PL.
Both significant wins again share the same (generator, ranker) cell,
Flan-T5-Small/BM25, on different tasks. RandMMR's own dial reaches further than the
single $\alpha{=}4$ point shown here: at two additional $\alpha$ values we ran for
LaMP-1 ($8$ and $20$), EE-D drops as low as 0.11--0.14 for the T5-Small cells ---
a fairness level plain PL only reaches at $\alpha{=}1$.

Read in aggregate, this is close to a wash: the win rate (43\%; mean
$\Delta$EU$_{\text{norm}}{=}{+}0.003$ across all 28 cells) sits near chance. What
is more robust is the asymmetry in the per-cell tests: not one of the 28 cells
reaches significance in PL's favor, while two do in RandMMR's --- both on the same
(generator, ranker) pair, Flan-T5-Small/BM25, at different tasks. We read this as
weak, cell-concentrated evidence for a utility edge rather than a broad one: real
enough that it never flips to a significant loss anywhere, but not something the
aggregate win rate supports on its own. The fairness side is unambiguous
regardless: RandMMR buys substantially lower EE-D than any workable plain-PL
setting except $\alpha{=}1$, without a significant utility cost anywhere in the
grid --- the fairness gain looks close to free, even though the utility picture
above it is closer to neutral than positive.

\FloatBarrier
\section{Conclusions}
\label{sec:conclusion}

\textbf{RQ1.} We find a narrow, asymmetric form of support for the
diversity-mediation hypothesis in our proposal, not the broad, uniform version it was
originally framed as: diversification helps mostly through MMR's relevance-aware
selection, and only in about half the tasks we test, not through PL's
fairness-motivated randomization on its own (Section~\ref{sec:rq1}).

\textbf{RQ2.} RandMMR, our proposed hybrid, is close to a wash on aggregate utility
against matched plain PL (12 of 28 tested cells, 43\%) but the two cells that do
reach significance both favor RandMMR and none favor PL, concentrated in a single
(generator, ranker) pair; it reliably reaches EE-D levels far below any workable
plain-PL setting except $\alpha{=}1$ (Section~\ref{sec:rq2}). The fairness gain
looks close to free; the utility picture on top of it is closer to neutral than
positive.

\section{Future Work}
\label{sec:future}

\begin{enumerate}[leftmargin=1.1em,itemsep=3pt]
  \item \textbf{Scale and realism.} Nearly every number in this paper comes
  from a setup chosen for what we could afford computationally, and as this is a course project, compromises were made:
  small sample counts, generators far below production scale, a ground
  truth that depends on those same small generators' own judgments, and a
  ROUGE-L-based utility signal standing in for a more precise but far costlier
  LLM-as-judge evaluation on the four generation tasks (Section~\ref{sec:threats}).
  Our own Flan-T5 models (77M/250M) are one to two
  orders of magnitude smaller than even the smallest generators commonly
  deployed in production; the source paper's largest, Flan-UL2 ($\sim$20B), is
  not undersized in the same way, but it is old --- the same 2022-era,
  non-RAG-finetuned family, several generations behind what is typically
  deployed today. A repeat at industry scale, with bigger sample counts, a
  dataset whose relevance labels do not depend on a single generator's own
  utility delta, and utility scored closer to human judgment than ROUGE-L
  allows, would show how much of RQ1--RQ2 survives outside this regime.
  \item \textbf{Task-specific RAG fairness-diversity balancing.} An interesting course of research in our opinion would be to study how much diversity is beneficial for different tasks or queries and how it can be balanced with fairness constraints in an optimal way.
  \item \textbf{Other reranking designs.} RandMMR is one specific point in a
  larger design space --- PL sampling for exactly the top rank, MMR for
  everything after it. Comparing it against other splits of that same idea,
  and against diversity/fairness methods outside the PL/MMR family entirely
  (e.g.\ xQuAD-style sub-topic coverage), could be an interesting research direction and may possibly yield another diversity-fairness balancing algorithm.
\end{enumerate}

{\footnotesize
\begin{thebibliography}{99}

\bibitem{eunkim2025fairrag}
T.\ Eun Kim and F.\ Diaz.
\newblock Towards Fair RAG: On the Impact of Fair Ranking in Retrieval-Augmented
Generation.
\newblock \emph{ICTIR}, 2025. arXiv:2409.11598.

\bibitem{castillo2018fairness}
C.\ Castillo.
\newblock Fairness and Transparency in Ranking.
\newblock \emph{ACM SIGIR Forum}, 52(2), 2018.

\bibitem{diaz2020evaluating}
F.\ Diaz, B.\ Mitra, M.\ D.\ Ekstrand, A.\ J.\ Biega, and B.\ Carterette.
\newblock Evaluating Stochastic Rankings with Expected Exposure.
\newblock \emph{CIKM}, 2020.

\bibitem{salemi2023lamp}
A.\ Salemi, S.\ Mysore, M.\ Bendersky, and H.\ Zamani.
\newblock LaMP: When Large Language Models Meet Personalization.
\newblock arXiv:2304.11406, 2023.

\bibitem{chung2024scaling}
H.\ W.\ Chung et al.
\newblock Scaling Instruction-Finetuned Language Models.
\newblock \emph{JMLR}, 25, 2024.

\bibitem{carbonell1998use}
J.\ Carbonell and J.\ Goldstein.
\newblock The Use of MMR, Diversity-Based Reranking for Reordering Documents and
Producing Summaries.
\newblock \emph{SIGIR}, 1998.

\bibitem{smyth2001similarity}
B.\ Smyth and P.\ McClave.
\newblock Similarity vs.\ Diversity.
\newblock \emph{ICCBR}, 2001.

\bibitem{plackett1975analysis}
R.\ L.\ Plackett.
\newblock The Analysis of Permutations.
\newblock \emph{Journal of the Royal Statistical Society: Series C}, 24(2), 1975.

\bibitem{luce1959individual}
R.\ D.\ Luce.
\newblock \emph{Individual Choice Behavior: A Theoretical Analysis}.
\newblock Wiley, 1959.

\end{thebibliography}
}

\appendix
\section{Deterministic-Baseline EU Detail}
\label{app:baseline}

Table~\ref{tab:baseline} compares our reproduction against the source paper's own
Table 2 two ways at once: a single deterministic-baseline EU$_{\text{norm}}$ value
per (retriever, generator) cell (Section~\ref{sec:method} for the formula) and,
following the source paper's own Table 2 exactly, the mean utility gain of
PL-sampled rankings over that baseline, pooled across all four PL $\alpha$ values
and binned into the same five EE-D disparity intervals the paper uses. Each row is
one (retriever, generator) cell for one of the two tasks the paper publishes this
comparison for; every cell shows Ours/Paper together.

On LaMP-1, a binary decision (accuracy over two candidate labels), the divergence is
generator-driven: Flan-T5-Base's baseline stays within 1--2\% of the paper, while
Flan-T5-Small diverges 46--51\% \emph{above} it. We traced this to a
content-invariant failure mode specific to Small on this task (Section~\ref{sec:method}):
it predicts the same label regardless of which document is retrieved, so its accuracy
reduces to how often that fixed label happens to be correct rather than reflecting
ranking quality --- and because gold\_ceiling is computed from that same
content-invariant behavior, it tracks the deterministic baseline almost exactly for
whichever queries the fixed label happens to get right, making the normalized value
unstable rather than a meaningful closeness-to-optimal measure. This makes Small's
LaMP-1 numbers intrinsically noisy at this sample size, independent of anything about
the reranker.

On LaMP-4, all four cells diverge in the same direction, above the paper, and now
much more uniformly than LaMP-1: 27--30\% for Flan-T5-Base, 28--37\% for
Flan-T5-Small --- the two generators sit far closer together here than on LaMP-1,
and Flan-T5-Small/BM25 no longer stands out as an outlier (37\%, in line with the
other three cells). This is a normalization-ceiling gap, not a generation-quality
one: our model\_ceiling is the best utility any rerank method we ran ever observed
for a query --- PL at any $\alpha$, MMR at any $\lambda$, or deterministic ---
taken together with gold\_ceiling above, rather than an exhaustive search over the
many permutations and subsets of a query's useful documents that could occupy the
$k{=}5$ slots. That search was not computationally feasible at this project's scale
(Section~\ref{sec:threats}), so our ceiling likely still sits below the paper's own
whenever a query has more useful documents than fit in $k{=}5$; since EU is raw
utility divided by this ceiling, a lower ceiling mechanically inflates our
normalized value for the same underlying generation quality.

The interval columns add a view the single baseline number cannot: how utility gain
over baseline changes as fairness (lower EE-D) increases. Six of the eight cells are
negative and significant at the lowest disparity interval ($[0,0.2)$, where the
fairness intervention is most aggressive), matching the paper's own pattern there;
three of those stay significant through the second interval before recovering
toward zero. Only Flan-T5-Small's two LaMP-1 cells show no significant interval in
either direction, consistent with their noisiness at this sample size (51 queries)
already discussed above. No cell now diverges sharply in \emph{shape} from the
paper's own pattern the way one did under the previous, PL-only ceiling:
Flan-T5-Small/BM25 on LaMP-4 was previously negative and significant across four of
its five intervals, unlike the paper's single significant interval for that cell;
it now matches the paper's one-interval shape, which we read as a sign that
broadening the ceiling, not any property specific to that cell, was driving the
earlier divergence.

The \emph{relative} pattern the paper's own Table 2 shows --- Flan-T5-Base above
Flan-T5-Small, BM25 above Contriever on LaMP-1, Contriever above BM25 on LaMP-4 ---
now survives almost entirely. The single exception is LaMP-4/BM25, where Small edges
narrowly ahead of Base (0.30 vs.\ 0.29); every other comparison, including both
retriever orderings for Flan-T5-Small (BM25 above Contriever on LaMP-1, Contriever
above BM25 on LaMP-4), now matches the paper. Under the previous PL-only ceiling,
Small violated both of those retriever orderings; broadening the ceiling resolved
them along with the interval-column outlier noted above.

\begin{table*}[t]
\centering
\scriptsize
\caption{Deterministic-baseline EU$_{\text{norm}}$ and EE-D-interval utility gain over
baseline: our reproduction vs.\ the source paper's own Table 2 (Section~\ref{sec:method}),
Ours/Paper in every cell. $^*p{<}0.05$ (ours: Welch's $t$-test vs.\ pooled baseline;
paper: unpaired Student's, per their Table 2 caption).}
\label{tab:baseline}
\begin{tabular}{@{}llcccccc@{}}
\toprule
Task & Retriever+Gen. & Baseline & [.0,.2) & [.2,.4) & [.4,.6) & [.6,.8) & [.8,1.0] \\
\midrule
LaMP-1 & BM25+Base   & 0.68/0.67 & $-$0.22$^{*}$/$-$0.20$^{*}$ & $-$0.10$^{*}$/$-$0.04 & $-$0.04/$-$0.08 & $-$0.05/$-$0.05 & $-$0.02/$-$0.02 \\
LaMP-1 & Contr.+Base & 0.64/0.64 & $-$0.14$^{*}$/$-$0.16$^{*}$ & $+$0.01/$+$0.05 & $-$0.06/$-$0.06 & $+$0.05/$+$0.03 & $+$0.01/$+$0.00 \\
LaMP-1 & BM25+Small  & 0.45/0.31 & $-$0.06/$-$0.12 & $-$0.42/$-$0.13 & $+$0.08/$-$0.18 & $+$0.05/$-$0.02 & $-$0.03/$-$0.15 \\
LaMP-1 & Contr.+Small& 0.43/0.29 & $-$0.03/$-$0.08 & $-$0.18/$-$0.29 & $-$0.18/$-$0.06 & $+$0.17/$+$0.03 & $-$0.03/$-$0.14 \\
\midrule
LaMP-4 & BM25+Base   & 0.29/0.22 & $-$0.08$^{*}$/$-$0.06$^{*}$ & $-$0.02/$+$0.00 & $+$0.02/$+$0.03 & $+$0.00/$+$0.01 & $+$0.00/$+$0.02 \\
LaMP-4 & Contr.+Base & 0.34/0.27 & $-$0.13$^{*}$/$-$0.10$^{*}$ & $-$0.04$^{*}$/$-$0.02 & $+$0.00/$+$0.01 & $-$0.01/$+$0.00 & $+$0.00/$+$0.01 \\
LaMP-4 & BM25+Small  & 0.30/0.22 & $-$0.11$^{*}$/$-$0.06$^{*}$ & $-$0.02/$+$0.00 & $-$0.01/$+$0.02 & $-$0.05/$+$0.01 & $-$0.00/$+$0.00 \\
LaMP-4 & Contr.+Small& 0.33/0.25 & $-$0.10$^{*}$/$-$0.09$^{*}$ & $-$0.04$^{*}$/$-$0.02$^{*}$ & $-$0.01/$+$0.00 & $+$0.02/$+$0.01 & $+$0.00/$+$0.00 \\
\bottomrule
\end{tabular}
\end{table*}

\section{Full RandMMR vs.\ Plain-PL Comparison}
\label{app:randmmr-full}

Table~\ref{tab:hybrid-full} gives the complete per-cell detail behind
Section~\ref{sec:rq2}'s 12-of-28 tally (Table~\ref{tab:hybrid} in the main text shows
only the two tasks with a published baseline). Columns are as in
Table~\ref{tab:hybrid}: PL $\alpha$ is the plain-PL condition matched to RandMMR's own
EE-D for that cell (uniformly $\alpha{=}2$ here), EE-D/EU pairs are macro means, and
$\Delta$EU is RandMMR's EU$_{\text{norm}}$ minus the matched PL's, both recomputed at
a uniform $N{=}10$ samples/query for every cell in the table.

\begin{table*}[t]
\centering
\scriptsize
\caption{RandMMR ($\alpha{=}4$, $\lambda\sim U(0.8,1.0)$) vs.\ matched plain PL, all
$7\times4{=}28$ cells, both conditions at $N{=}10$. Two of the 28 cells reach
significance, both favoring RandMMR (Flan-T5-Small/BM25 on LaMP-4, $p{<}0.01$, and
on LaMP-2, $p{<}0.05$); Welch's $t$-test vs.\ matched PL.}
\label{tab:hybrid-full}
\begin{tabular}{@{}lllrllr@{}}
\toprule
Task & Generator & Ranker & PL $\alpha$ & PL EE-D/EU & RandMMR EE-D/EU & $\Delta$EU \\
\midrule
\multirow{4}{*}{LaMP-1} & T5-Base & BM25 & 2 & 0.19/0.51 & 0.19/0.49 & $-0.019$ \\
 & T5-Base & Contriever & 2 & 0.18/0.56 & 0.17/0.51 & $-0.048$ \\
 & T5-Small & BM25 & 2 & 0.18/0.40 & 0.17/0.39 & $-0.008$ \\
 & T5-Small & Contriever & 2 & 0.17/0.39 & 0.16/0.41 & $+0.020$ \\
\cmidrule(lr){1-7}
\multirow{4}{*}{LaMP-2} & T5-Base & BM25 & 2 & 0.38/0.32 & 0.36/0.31 & $-0.008$ \\
 & T5-Base & Contriever & 2 & 0.38/0.32 & 0.36/0.31 & $-0.006$ \\
 & T5-Small & BM25 & 2 & 0.34/0.24 & 0.34/0.35 & $+0.108^{*}$ \\
 & T5-Small & Contriever & 2 & 0.36/0.38 & 0.34/0.37 & $-0.011$ \\
\cmidrule(lr){1-7}
\multirow{4}{*}{LaMP-3} & T5-Base & BM25 & 2 & 0.15/0.93 & 0.15/0.93 & $-0.001$ \\
 & T5-Base & Contriever & 2 & 0.15/0.92 & 0.14/0.92 & $-0.001$ \\
 & T5-Small & BM25 & 2 & 0.14/0.87 & 0.14/0.89 & $+0.016$ \\
 & T5-Small & Contriever & 2 & 0.15/0.89 & 0.14/0.89 & $-0.000$ \\
\cmidrule(lr){1-7}
\multirow{4}{*}{LaMP-4} & T5-Base & BM25 & 2 & 0.25/0.25 & 0.25/0.25 & $+0.003$ \\
 & T5-Base & Contriever & 2 & 0.24/0.25 & 0.24/0.25 & $+0.000$ \\
 & T5-Small & BM25 & 2 & 0.26/0.22 & 0.25/0.27 & $+0.050^{**}$ \\
 & T5-Small & Contriever & 2 & 0.24/0.26 & 0.23/0.26 & $-0.007$ \\
\cmidrule(lr){1-7}
\multirow{4}{*}{LaMP-5} & T5-Base & BM25 & 2 & 0.18/0.53 & 0.17/0.53 & $+0.002$ \\
 & T5-Base & Contriever & 2 & 0.18/0.53 & 0.18/0.53 & $+0.002$ \\
 & T5-Small & BM25 & 2 & 0.18/0.42 & 0.17/0.41 & $-0.007$ \\
 & T5-Small & Contriever & 2 & 0.18/0.41 & 0.18/0.41 & $-0.006$ \\
\cmidrule(lr){1-7}
\multirow{4}{*}{LaMP-6} & T5-Base & BM25 & 2 & 0.24/0.41 & 0.23/0.41 & $+0.001$ \\
 & T5-Base & Contriever & 2 & 0.24/0.40 & 0.23/0.40 & $-0.001$ \\
 & T5-Small & BM25 & 2 & 0.24/0.39 & 0.23/0.39 & $-0.001$ \\
 & T5-Small & Contriever & 2 & 0.24/0.38 & 0.23/0.38 & $-0.002$ \\
\cmidrule(lr){1-7}
\multirow{4}{*}{LaMP-7} & T5-Base & BM25 & 2 & 0.50/0.66 & 0.50/0.65 & $-0.007$ \\
 & T5-Base & Contriever & 2 & 0.49/0.65 & 0.46/0.65 & $+0.004$ \\
 & T5-Small & BM25 & 2 & 0.54/0.53 & 0.53/0.54 & $+0.009$ \\
 & T5-Small & Contriever & 2 & 0.52/0.52 & 0.50/0.53 & $+0.001$ \\
\bottomrule
\end{tabular}
\end{table*}

\section{Full ILD-Quartile Utility Table}
\label{app:ild-quartiles-full}

Table~\ref{tab:ild-quartiles-full} gives the complete per-cell detail behind
Table~\ref{tab:ild-quartiles} (Section~\ref{sec:rq1}): the same ILD$_{\text{norm}}$
quartile-binned utility gain over deterministic, for all seven LaMP tasks rather
than only the two with a published baseline. Two cells (Flan-T5-Base/Contriever and
Flan-T5-Small/Contriever, both MMR on LaMP-2) have too few distinct
ILD$_{\text{norm}}$ values to split into four quartiles and are marked n/a.

\begin{table*}[t]
\centering
\scriptsize
\caption{PL and MMR utility gain over deterministic, by each generator/ranker/method
cell's own ILD$_{\text{norm}}$ quartile, all seven LaMP tasks. Bracketed values are
the quartile's own ILD$_{\text{norm}}$ boundaries. $^*p{<}0.05$, $^{**}p{<}0.01$
(Welch's $t$-test vs.\ deterministic).}
\label{tab:ild-quartiles-full}
\begin{tabular}{@{}lllrlrrrr@{}}
\toprule
Task & Gen. & Ranker & Base & Method & Q1 (least diverse) & Q2 & Q3 & Q4 (most diverse) \\
\midrule
\multirow{2}{*}{LaMP-1} & \multirow{2}{*}{T5-Base} & \multirow{2}{*}{BM25} & \multirow{2}{*}{0.681} & PL & [0.73,0.97] $-$0.071 & [0.97,0.98] $-$0.083$^{*}$ & [0.98,0.99] $-$0.121$^{**}$ & [0.99,1.00] $-$0.165$^{**}$ \\
 & & & & MMR & [0.73,0.97] $+$0.053 & [0.97,0.98] $-$0.039 & [0.98,0.99] $+$0.002 & [0.99,1.00] $-$0.089$^{*}$ \\
\cmidrule(lr){1-9}
 & \multirow{2}{*}{T5-Base} & \multirow{2}{*}{Contr.} & \multirow{2}{*}{0.642} & PL & [0.81,0.96] $-$0.042 & [0.96,0.98] $-$0.061 & [0.98,0.99] $-$0.021 & [0.99,1.00] $-$0.081$^{*}$ \\
 & & & & MMR & [0.85,0.97] $-$0.007 & [0.97,0.98] $+$0.020 & [0.98,0.99] $+$0.067 & [0.99,1.00] $-$0.094$^{*}$ \\
\cmidrule(lr){1-9}
 & \multirow{2}{*}{T5-Small} & \multirow{2}{*}{BM25} & \multirow{2}{*}{0.451} & PL & [0.75,0.96] $-$0.135 & [0.96,0.98] $+$0.051 & [0.98,0.99] $+$0.047 & [0.99,1.00] $-$0.071 \\
 & & & & MMR & [0.90,0.96] $-$0.143 & [0.96,0.98] $+$0.089 & [0.98,1.00] $+$0.002 & [1.00,1.00] $-$0.054 \\
\cmidrule(lr){1-9}
 & \multirow{2}{*}{T5-Small} & \multirow{2}{*}{Contr.} & \multirow{2}{*}{0.431} & PL & [0.81,0.96] $-$0.043 & [0.96,0.97] $-$0.047 & [0.97,0.99] $-$0.022 & [0.99,1.00] $+$0.035 \\
 & & & & MMR & [0.87,0.96] $-$0.041 & [0.96,0.98] $+$0.031 & [0.98,0.99] $+$0.019 & [0.99,1.00] $+$0.022 \\
\midrule
\multirow{2}{*}{LaMP-2} & \multirow{2}{*}{T5-Base} & \multirow{2}{*}{BM25} & \multirow{2}{*}{0.300} & PL & [0.92,0.97] $+$0.021 & [0.97,0.98] $-$0.005 & [0.98,0.99] $-$0.021 & [0.99,1.00] $+$0.081$^{*}$ \\
 & & & & MMR & [0.92,0.97] $+$0.032 & [0.97,0.98] $-$0.016 & [0.98,0.99] $-$0.064 & [0.99,1.00] $+$0.062 \\
\cmidrule(lr){1-9}
 & \multirow{2}{*}{T5-Base} & \multirow{2}{*}{Contr.} & \multirow{2}{*}{0.307} & PL & [0.95,0.98] $+$0.005 & [0.98,0.99] $+$0.012 & [0.99,1.00] $+$0.055 & [1.00,1.00] $-$0.004 \\
 & & & & MMR & n/a & n/a & n/a & n/a \\
\cmidrule(lr){1-9}
 & \multirow{2}{*}{T5-Small} & \multirow{2}{*}{BM25} & \multirow{2}{*}{0.333} & PL & [0.94,0.98] $-$0.057 & [0.98,0.99] $+$0.011 & [0.99,0.99] $-$0.060 & [0.99,1.00] $+$0.128$^{**}$ \\
 & & & & MMR & [0.94,0.97] $-$0.177$^{**}$ & [0.97,0.98] $+$0.036 & [0.98,0.99] $+$0.040 & [0.99,1.00] $+$0.071 \\
\cmidrule(lr){1-9}
 & \multirow{2}{*}{T5-Small} & \multirow{2}{*}{Contr.} & \multirow{2}{*}{0.375} & PL & [0.95,0.98] $+$0.016 & [0.98,0.99] $-$0.041 & [0.99,1.00] $-$0.016 & [1.00,1.00] $+$0.032 \\
 & & & & MMR & n/a & n/a & n/a & n/a \\
\midrule
\multirow{2}{*}{LaMP-3} & \multirow{2}{*}{T5-Base} & \multirow{2}{*}{BM25} & \multirow{2}{*}{0.923} & PL & [0.00,0.94] $+$0.005 & [0.94,0.97] $+$0.003 & [0.97,0.98] $-$0.003 & [0.98,1.00] $+$0.001 \\
 & & & & MMR & [0.00,0.95] $-$0.003 & [0.95,0.97] $+$0.020$^{*}$ & [0.97,0.98] $-$0.012 & [0.98,1.00] $-$0.005 \\
\cmidrule(lr){1-9}
 & \multirow{2}{*}{T5-Base} & \multirow{2}{*}{Contr.} & \multirow{2}{*}{0.927} & PL & [0.00,0.95] $-$0.004 & [0.95,0.97] $+$0.003 & [0.97,0.99] $-$0.002 & [0.99,1.00] $+$0.001 \\
 & & & & MMR & [0.00,0.96] $+$0.009 & [0.96,0.97] $+$0.004 & [0.97,0.99] $+$0.000 & [0.99,1.00] $+$0.003 \\
\cmidrule(lr){1-9}
 & \multirow{2}{*}{T5-Small} & \multirow{2}{*}{BM25} & \multirow{2}{*}{0.908} & PL & [0.00,0.95] $-$0.006 & [0.95,0.98] $-$0.010 & [0.98,0.99] $-$0.019 & [0.99,1.00] $-$0.023$^{*}$ \\
 & & & & MMR & [0.00,0.95] $-$0.015 & [0.95,0.97] $-$0.016 & [0.97,0.98] $-$0.007 & [0.98,1.00] $-$0.008 \\
\cmidrule(lr){1-9}
 & \multirow{2}{*}{T5-Small} & \multirow{2}{*}{Contr.} & \multirow{2}{*}{0.897} & PL & [0.00,0.96] $+$0.010 & [0.96,0.98] $-$0.004 & [0.98,0.99] $-$0.012 & [0.99,1.00] $-$0.004 \\
 & & & & MMR & [0.00,0.96] $+$0.019 & [0.96,0.98] $+$0.003 & [0.98,0.99] $-$0.021 & [0.99,1.00] $-$0.012 \\
\midrule
\multirow{2}{*}{LaMP-4} & \multirow{2}{*}{T5-Base} & \multirow{2}{*}{BM25} & \multirow{2}{*}{0.291} & PL & [0.46,0.97] $-$0.011 & [0.97,0.98] $-$0.028$^{*}$ & [0.98,0.99] $-$0.035$^{**}$ & [0.99,1.00] $-$0.022 \\
 & & & & MMR & [0.46,0.97] $+$0.003 & [0.97,0.98] $+$0.021 & [0.98,0.99] $-$0.006 & [0.99,1.00] $-$0.005 \\
\cmidrule(lr){1-9}
 & \multirow{2}{*}{T5-Base} & \multirow{2}{*}{Contr.} & \multirow{2}{*}{0.340} & PL & [0.50,0.97] $-$0.002 & [0.97,0.98] $-$0.060$^{**}$ & [0.98,0.99] $-$0.057$^{**}$ & [0.99,1.00] $-$0.071$^{**}$ \\
 & & & & MMR & [0.52,0.97] $+$0.029$^{*}$ & [0.97,0.99] $-$0.005 & [0.99,1.00] $-$0.033$^{*}$ & [1.00,1.00] $-$0.005 \\
\cmidrule(lr){1-9}
 & \multirow{2}{*}{T5-Small} & \multirow{2}{*}{BM25} & \multirow{2}{*}{0.297} & PL & [0.47,0.97] $-$0.010 & [0.97,0.98] $-$0.021 & [0.98,0.99] $-$0.005 & [0.99,1.00] $-$0.001 \\
 & & & & MMR & [0.47,0.97] $+$0.006 & [0.97,0.98] $-$0.002 & [0.98,0.99] $-$0.009 & [0.99,1.00] $+$0.019 \\
\cmidrule(lr){1-9}
 & \multirow{2}{*}{T5-Small} & \multirow{2}{*}{Contr.} & \multirow{2}{*}{0.326} & PL & [0.51,0.98] $-$0.004 & [0.98,0.99] $-$0.059$^{**}$ & [0.99,0.99] $-$0.033$^{**}$ & [0.99,1.00] $-$0.052$^{**}$ \\
 & & & & MMR & [0.53,0.97] $+$0.028 & [0.97,0.99] $-$0.012 & [0.99,1.00] $+$0.011 & [1.00,1.00] $-$0.047$^{**}$ \\
\midrule
\multirow{2}{*}{LaMP-5} & \multirow{2}{*}{T5-Base} & \multirow{2}{*}{BM25} & \multirow{2}{*}{0.531} & PL & [0.42,0.96] $-$0.021 & [0.96,0.98] $+$0.003 & [0.98,0.99] $+$0.003 & [0.99,1.00] $+$0.018 \\
 & & & & MMR & [0.42,0.96] $-$0.034$^{*}$ & [0.96,0.98] $+$0.020 & [0.98,0.99] $+$0.018 & [0.99,1.00] $+$0.047$^{**}$ \\
\cmidrule(lr){1-9}
 & \multirow{2}{*}{T5-Base} & \multirow{2}{*}{Contr.} & \multirow{2}{*}{0.550} & PL & [0.42,0.96] $-$0.012 & [0.96,0.98] $-$0.021 & [0.98,0.99] $-$0.016 & [0.99,1.00] $-$0.007 \\
 & & & & MMR & [0.42,0.95] $-$0.034$^{*}$ & [0.95,0.97] $-$0.021 & [0.97,0.99] $+$0.017 & [0.99,1.00] $+$0.042$^{*}$ \\
\cmidrule(lr){1-9}
 & \multirow{2}{*}{T5-Small} & \multirow{2}{*}{BM25} & \multirow{2}{*}{0.458} & PL & [0.43,0.97] $-$0.008 & [0.97,0.99] $-$0.009 & [0.99,1.00] $-$0.022 & [1.00,1.00] $-$0.045$^{**}$ \\
 & & & & MMR & [0.43,0.96] $+$0.028 & [0.96,0.98] $-$0.036 & [0.98,0.99] $+$0.016 & [0.99,1.00] $+$0.026 \\
\cmidrule(lr){1-9}
 & \multirow{2}{*}{T5-Small} & \multirow{2}{*}{Contr.} & \multirow{2}{*}{0.465} & PL & [0.43,0.96] $-$0.004 & [0.96,0.98] $-$0.040$^{*}$ & [0.98,0.99] $-$0.025 & [0.99,1.00] $-$0.054$^{**}$ \\
 & & & & MMR & [0.43,0.95] $+$0.008 & [0.95,0.97] $-$0.033 & [0.97,0.99] $-$0.031 & [0.99,1.00] $+$0.020 \\
\midrule
\multirow{2}{*}{LaMP-6} & \multirow{2}{*}{T5-Base} & \multirow{2}{*}{BM25} & \multirow{2}{*}{0.473} & PL & [0.02,0.78] $-$0.028 & [0.78,0.92] $-$0.056$^{**}$ & [0.92,0.97] $-$0.036$^{*}$ & [0.97,1.00] $-$0.048$^{**}$ \\
 & & & & MMR & [0.02,0.82] $+$0.040$^{*}$ & [0.82,0.94] $-$0.014 & [0.94,0.98] $-$0.017 & [0.98,1.00] $-$0.034 \\
\cmidrule(lr){1-9}
 & \multirow{2}{*}{T5-Base} & \multirow{2}{*}{Contr.} & \multirow{2}{*}{0.444} & PL & [0.02,0.78] $-$0.003 & [0.78,0.91] $-$0.059$^{**}$ & [0.91,0.97] $-$0.034 & [0.97,1.00] $+$0.001 \\
 & & & & MMR & [0.02,0.85] $+$0.023 & [0.85,0.94] $-$0.041$^{*}$ & [0.94,0.98] $-$0.023 & [0.98,1.00] $+$0.012 \\
\cmidrule(lr){1-9}
 & \multirow{2}{*}{T5-Small} & \multirow{2}{*}{BM25} & \multirow{2}{*}{0.442} & PL & [0.02,0.83] $-$0.035 & [0.83,0.95] $-$0.047$^{*}$ & [0.95,0.98] $-$0.024 & [0.98,1.00] $-$0.030 \\
 & & & & MMR & [0.02,0.85] $-$0.002 & [0.85,0.95] $+$0.009 & [0.95,0.98] $-$0.017 & [0.98,1.00] $-$0.027 \\
\cmidrule(lr){1-9}
 & \multirow{2}{*}{T5-Small} & \multirow{2}{*}{Contr.} & \multirow{2}{*}{0.424} & PL & [0.02,0.83] $-$0.023 & [0.83,0.94] $-$0.074$^{**}$ & [0.94,0.98] $-$0.030 & [0.98,1.00] $+$0.032 \\
 & & & & MMR & [0.02,0.88] $+$0.007 & [0.88,0.95] $-$0.027 & [0.95,0.98] $-$0.018 & [0.98,1.00] $+$0.075$^{**}$ \\
\midrule
\multirow{2}{*}{LaMP-7} & \multirow{2}{*}{T5-Base} & \multirow{2}{*}{BM25} & \multirow{2}{*}{0.675} & PL & [0.00,0.97] $-$0.032 & [0.97,0.98] $-$0.002 & [0.98,0.99] $+$0.019 & [0.99,1.00] $-$0.035 \\
 & & & & MMR & [0.00,0.97] $-$0.068$^{*}$ & [0.97,0.99] $+$0.083$^{**}$ & [0.99,1.00] $+$0.003 & [1.00,1.00] $-$0.083$^{**}$ \\
\cmidrule(lr){1-9}
 & \multirow{2}{*}{T5-Base} & \multirow{2}{*}{Contr.} & \multirow{2}{*}{0.645} & PL & [0.00,0.97] $-$0.001 & [0.97,0.98] $-$0.007 & [0.98,0.99] $+$0.020 & [0.99,1.00] $+$0.006 \\
 & & & & MMR & [0.00,0.98] $-$0.001 & [0.98,0.99] $-$0.043 & [0.99,1.00] $-$0.007 & [1.00,1.00] $+$0.035 \\
\cmidrule(lr){1-9}
 & \multirow{2}{*}{T5-Small} & \multirow{2}{*}{BM25} & \multirow{2}{*}{0.535} & PL & [0.00,0.98] $+$0.030 & [0.98,0.99] $-$0.009 & [0.99,0.99] $-$0.031 & [0.99,1.00] $+$0.012 \\
 & & & & MMR & [0.00,0.97] $+$0.033 & [0.97,0.99] $-$0.034 & [0.99,1.00] $-$0.029 & [1.00,1.00] $-$0.042 \\
\cmidrule(lr){1-9}
 & \multirow{2}{*}{T5-Small} & \multirow{2}{*}{Contr.} & \multirow{2}{*}{0.530} & PL & [0.00,0.98] $-$0.000 & [0.98,0.99] $-$0.025 & [0.99,1.00] $-$0.004 & [1.00,1.00] $+$0.021 \\
 & & & & MMR & [0.00,0.98] $-$0.017 & [0.98,0.99] $-$0.055$^{*}$ & [0.99,1.00] $+$0.003 & [1.00,1.00] $+$0.002 \\
\bottomrule
\end{tabular}
\end{table*}

\end{document}
